% Save in sections/ as section-01-name-content.tex.
% This fragment is shared by the chapter and assembled course notes.
% Do not compile directly.
\section{First Section Title}

Every chapter builds upon rigorous narrative exposition. Modern reading materials must bridge the gap between theoretical principles and practical engineering execution.

\begin{DefinitionBox}{Core Mathematical Entity}
  Let \(\mathcal{X}\) denote the sample space, and let \(\theta \in \Theta\) represent the unobserved parameter vector. The probability model is parameterized by density \(p(x \mid \theta)\).
\end{DefinitionBox}

\subsection{Analytical Derivation and Mechanics}

To understand how parameters behave, consider the integral transformation across the measurable domain:

\BookEquation{p(\theta \mid y) = \frac{p(y \mid \theta) \, p(\theta)}{\int_\Theta p(y \mid \theta') \, p(\theta') \, d\theta'}}

\begin{InsightBox}{Interpretation of the Normalizer}
  The evidence integral in the denominator ensures that the posterior distribution sums or integrates strictly to unity, making it a valid probability distribution rather than merely a relative likelihood profile.
\end{InsightBox}

\begin{ExampleBox}{One-Parameter Normal Model}
  Suppose observations \(y_1, \dots, y_n \sim \mathcal{N}(\mu, \sigma^2)\) with known variance \(\sigma^2\) and a conjugate normal prior \(\mu \sim \mathcal{N}(\mu_0, \sigma_0^2)\). The resulting posterior mean satisfies:
  \[
    \mu_n = \frac{\frac{1}{\sigma_0^2}\mu_0 + \frac{n}{\sigma^2}\bar{y}}{\frac{1}{\sigma_0^2} + \frac{n}{\sigma^2}}
  \]
  This is a precision-weighted convex combination of prior belief and empirical observation.
\end{ExampleBox}

\begin{WarningBox}{Beware of Improper Posteriors}
  Using flat or uninformative priors without verifying integrability can lead to improper posterior distributions where inference algorithms fail silently.
\end{WarningBox}
